% Generated draft: requires mathematical and source review.
\documentclass[11pt]{article}
\usepackage[margin=25mm]{geometry}
\usepackage{fontspec}
\setmainfont{Latin Modern Roman}
\usepackage{amsmath,amssymb,mathtools,unicode-math}
\setmathfont{Latin Modern Math}
\usepackage{enumitem,needspace,xurl,hyperref}
\setlength{\parindent}{0pt}
\setlength{\parskip}{5pt}
\setlength{\emergencystretch}{3em}
\newcommand{\N}{\mathbb{N}}
\newcommand{\Z}{\mathbb{Z}}
\newcommand{\Q}{\mathbb{Q}}
\newcommand{\R}{\mathbb{R}}
\newcommand{\C}{\mathbb{C}}
\newcommand{\F}{\mathbb{F}}
\newcommand{\A}{\mathbb{A}}
\newcommand{\PP}{\mathbb{P}}
\newcommand{\E}{\mathbb{E}}
\newcommand{\eps}{\varepsilon}
\newcommand{\dd}{\,\mathrm{d}}
\begin{document}
\section*{UnsolvedMath ID: 20002317}
\textbf{Source problem code:} catalogue entry
\subsection{Definitions and mathematical statement}
Conjecture 15 says that the asymptotically almost sure theory determined by a preceding \(0\)--\(1\) law is that of \(F_2\). The preceding random model is absent, and the extraction includes a page header and the next ``Standard systems of RCF (Marker)'' section. In the aid's recovered setting \(F_2\) is the rank-two free group, not the two-element field; the Marker heading does not attribute this conjecture to him.

Fix \(n\ge2\), write \(F_n=\langle x_1,\ldots,x_n\rangle\), and for a relator \(r\) set
\[
 G_r=F_n/\langle\!\langle r\rangle\!\rangle.
\]
The aid's resolved presentation model samples \(r\) uniformly among reduced relators of exact length \(L\), with the stated cyclic-reduction variants. It attributes to Kharlampovich--Miasnikov--Sklinos III at density \(d=0\) the pure-group-language limit
\[
 \lim_{L\to\infty}\Pr(G_r\models\sigma)=
 \begin{cases}1,&F_2\models\sigma,\\0,&F_2\not\models\sigma,\end{cases}
\]
for every first-order group sentence \(\sigma\). The target is independent of the fixed rank because \(F_n\equiv F_2\) for \(n\ge2\). Constants naming the chosen generators or extra predicates are not included in this claim. The exact theorem citation and all density-model conventions are not supplied.

The aid reports normalization robustness for any event depending only on the quotient group. If its probability has an endpoint limit \(0\) or \(1\), the same endpoint holds for uniform freely reduced spheres, uniform cyclically reduced spheres, and uniform unoriented cyclic-relator orbits. The last model quotients cyclic words by rotation and inversion before choosing each orbit with equal weight. For that cyclic model, pushing uniform word measure to orbits and taking uniform orbit measure have total-variation distance \(o(1)\) as \(L\to\infty\). This is not a total-variation claim identifying every freely reduced and cyclically reduced exact-length distribution, nor a statement about arbitrary intermediate limiting values.

Exponential sphere-to-ball averaging also gives the presentation-weighted version with relator length at most \(L\). These are all measures on presentations or relator orbits pushed to groups; different presentations yielding the same abstract group still contribute their sampling weights. The literal alternative of assigning equal probability to each distinct group isomorphism type is not established by these results and remains unresolved in the aid's assessment. Thus the reported presentation-model resolution is not silently substituted for an unspecified uniform-isomorphism-type reading of the missing preceding conjecture.
\subsection{Short English statement}
The contaminated excerpt predicts the rank-two free group's theory as a random-group limit. The aid reports the result for presentation-weighted one-relator models and shows endpoint robustness under common relator normalizations. Uniform weighting of distinct isomorphism types is a separate unresolved interpretation.
\subsection{Sources}
\begin{itemize}
\item Local UnsolvedMath registry, permanent ID 20002317, exact source code catalogue entry; canonical \texttt{aim-logic-notes.json}, \texttt{notes[92]}, tag \texttt{conjecture}; AIM Workshop Problem Lists, Logic, \emph{Computable stability theory}, Conjecture 15. Created and updated 14 August 2026.
\item Source URL \url{https://aimath.org/pastworkshops/computestabproblems.pdf}. The supplied verification identifies page-break contamination; the preceding \(0\)--\(1\)-law model is not reproduced. The PDF was not accessed.
\item Kharlampovich--Miasnikov--Sklinos III, density \(d=0\), is named for the exact-length presentation limit. No full title, date, identifier, theorem number, or separate URL is supplied.
\item Attempt 1: \texttt{partial\_result}, problem status at run \texttt{partially\_solved}, rigor label \texttt{proved}, novelty confidence \texttt{low}; classification \texttt{PARTIAL-PROGRESS}, suggested difficulty L4, registry status \texttt{partially\_solved}. Candidate: normalization robustness and cyclic-orbit total-variation comparison. The isomorphism-type qualification is part of the supplied assessment.
\item Ulam AI / UnsolvedMath Contributors; CC BY 4.0. Machine-generated research aid, not peer reviewed; independently verify before citation.
\end{itemize}
\end{document}
